% DEGraph -- arXiv preprint
% Build: compiles on Overleaf (pdfLaTeX). No local TeX toolchain was available
% when this was authored, so it has NOT been locally compiled -- do a first
% Overleaf build and fix any stragglers. Uses self-contained thebibliography
% (no external .bib / bibtex pass needed).
\documentclass[11pt]{article}

\usepackage[utf8]{inputenc}
\usepackage[T1]{fontenc}
\usepackage[margin=1in]{geometry}
\usepackage{microtype}
\usepackage{booktabs}
\usepackage{array}
\usepackage{amsmath,amssymb}
\usepackage{graphicx}
\usepackage{xcolor}
\usepackage{listings}
\usepackage[hidelinks]{hyperref}
\usepackage{underscore} % lets _ print literally in text (column names are everywhere)

\lstset{
  basicstyle=\ttfamily\small,
  breaklines=true,
  columns=fullflexible,
  keepspaces=true,
  frame=single,
  framesep=4pt,
  xleftmargin=4pt,
}

\newcommand{\code}[1]{\texttt{#1}}

% Ragged-right fixed-width column: wraps long \texttt content at spaces instead of
% overflowing the margin, and avoids the stretched inter-word spacing that
% justified narrow p-columns produce.
\newcolumntype{L}[1]{>{\raggedright\arraybackslash}p{#1}}

% Long \texttt spans cannot hyphenate, so give TeX slack to avoid overfull
% lines around short unbreakable tokens (# COMMAND, sqlglot, .ipynb).
\setlength{\emergencystretch}{3em}

\title{\textbf{What Breaks Static Column-Level Lineage in\\
Production PySpark? A Tool and Two Industrial Case Studies}}

\author{Rojan Raj Thapa\\ \texttt{rojaneverest@gmail.com}}
\date{\today}

\begin{document}
\maketitle

\begin{abstract}
Data engineers routinely need to answer change-impact questions about PySpark
pipelines before deploying a change: if a column is renamed or removed, which
downstream tables and columns break? Which source column feeds a given gold-layer
metric? Today these questions are answered either by reading source across many
files by hand, or by runtime lineage tools (OpenLineage, Spline, Unity Catalog)
that require the pipeline to actually execute---ruling out code under review,
experimental branches, and pre-merge CI checks. No existing static analyzer
recovers column-level data lineage from the PySpark DataFrame API: the closest
tool, \code{pyspark-ast-lineage}, extracts only a flat set of source and sink
\emph{table names}, with no intermediate DataFrames, no transformation edges, and
no column-level provenance.

We present DEGraph, a static analyzer that traverses PySpark repositories via
Python AST analysis and \code{sqlglot}---with no code execution and no LLM
inference at extraction time---and produces a typed, column-level data-lineage
graph: global nodes are persisted tables, file-local nodes are intermediate
DataFrames, and eight edge types (Reads, Writes, Derives, Filters, Projects,
Joins, Aggregates, OpaqueTransform) carry column-level provenance, merge keys,
conditional classification rules, and array-struct field provenance. On a
hand-labeled benchmark (\code{repo_synthetic_small}, 13 files, 8 tables, 57
ground-truth edges), the extractor recovers lineage at \textbf{98\%
precision/recall/F1} under node-name-agnostic (semantic) matching and 79\% under
exact intermediate-name matching---the gap being a labeling convention, not a
missed relationship. Against \code{pyspark-ast-lineage} on the same benchmark,
DEGraph recovers all 8 tables and 57 typed edges with column provenance, while the
baseline---though functional on its own example---recovers nothing, because real
Databricks idioms (f-string fully-qualified names resolved from \code{\%run}
config, streaming reads) defeat its variable resolver.

On top of this graph we build two execution-free capabilities that neither runtime
tools (which require execution) nor lexical retrieval (which has no structure) can
provide. \textbf{Static change-impact analysis} computes the forward transitive
closure over inverted column provenance: given a changed column, it returns every
downstream column and table affected. On 21 ground-truth change scenarios---spanning
the synthetic benchmarks \emph{and} real third-party Databricks Declarative-Pipeline
code---it achieves \textbf{100\% precision (zero false positives), 78\% recall, F1
88\%}. We state the precision property conditionally, and explain why: because impact
is reachability over edges the extractor actually emitted, DEGraph reports no false
positives \emph{when it has resolved the seed column's lineage}, and it is now built
to report nothing rather than guess when it has not.

We then evaluate on two proprietary production pipelines at an industrial partner, one used for
development and one \textbf{held out} behind a frozen tool version. The results are
mixed and we report them in full. Reaching usable coverage on the development
pipeline required handling \textbf{nine concrete syntactic idioms}---custom sink
classes, dict-returning table loaders, config-driven table names, bulk renames,
splatted and non-splatted projections---none of which was a deficiency of the graph
model; we argue the barrier to static column lineage in production is \emph{surface-form
coverage, not transformation semantics}. On that pipeline impact reached 100\%
precision at 68\% recall. On the held-out pipeline it reached \textbf{0.167 precision
and 0.118 recall}: the generalization claim, tested properly, does not hold, and
coverage is strongly domain-dependent. The held-out run also exposed a precision
defect invisible to every prior evaluation---an unresolved seed silently matched a
column name in a \emph{different} table---which we diagnose, repair, and verify
changes none of the previously reported results.
\textbf{Lineage version-diff} compares two source revisions and flags breaking
changes with their blast radius; on six realistic edits it detects added/removed
columns at 100\% precision/recall and classifies all three breaking and all three
safe edits correctly (zero false alarms), statically, before any code runs. As a
reliability check on the single-author ground truth, two independent LLM judges
re-derived the impact labels from source alone, agreeing at Cohen's $\kappa=0.86$
and $0.91$.

We also report an honest negative result that reshaped the project. We tested
whether serving the graph as a compact serialization is a superior \emph{LLM
context} format, against raw source and a token-matched BM25 retrieval baseline,
across four model families. At equal token budget, \textbf{BM25 matched or beat the
graph serialization} (pooled Gemini: raw 85.8\%, BM25 81.7\%, graph 75.0\%), so we
make \emph{no} claim that the serialization is a better LLM context layer. DEGraph's
contribution is instead the deterministic, static, execution-free structural
analysis (impact and diff) that retrieval and runtime tools cannot offer.
\end{abstract}

\section{Introduction}

\subsection{The problem}
Modern data engineering on cloud lakehouse platforms---Databricks, Snowflake,
Microsoft Fabric---organizes analytical data into a three-tier medallion
architecture: \emph{bronze} tables hold raw ingested data; \emph{silver} tables
hold cleaned, conformed, joined datasets derived from bronze; \emph{gold} tables
hold business-ready aggregates derived from silver. In production this is
implemented across hundreds of PySpark and SQL files: dozens of bronze ingestion
scripts, silver pipelines heavy with joins and type coercion, and gold aggregation
layers, mixing PySpark DataFrame calls, embedded \code{spark.sql()} strings,
standalone \code{.sql} files, and Databricks notebooks.

This scale creates a practical problem for \emph{change-impact analysis}. Before
merging a change, an engineer must answer questions like \emph{``if I rename
\code{lifetime_revenue} in \code{silver.\allowbreak customer_profile}, which downstream gold
tables and columns break?''} Answering correctly requires tracing data flow across
many files at column-level granularity. Two options exist today, both structurally
limited. The first is to read the source by hand---accurate but slow and
error-prone at scale. The second is a runtime lineage tool (Apache Atlas,
OpenLineage, Spline, Unity Catalog Lineage), but all of these \emph{observe}
lineage from an executing pipeline; they cannot analyze code that has not yet
run---a pull request under review, an experimental branch, or a pre-merge CI check.
The pre-deployment moment, when impact analysis is most valuable, is exactly when
runtime tools are unavailable.

The root cause is that the column-level data-flow structure these questions depend
on is never materialized statically. A static analyzer that recovers, from source
alone, \emph{which columns from which source tables flow into which derived columns
across which files} would make change-impact analysis a deterministic
query---answerable before any code runs, with no cluster and no execution. The
closest existing static tool, \code{pyspark-ast-lineage}, recovers only the
\emph{names} of source and sink tables; the transformation structure
between them is left opaque.

This paper therefore asks two questions rather than one. The first is a
\emph{capability} question---can such a graph be recovered from PySpark source, and
is it accurate enough to answer change-impact and version-diff queries? The second,
which we consider the more useful contribution, is an \emph{applicability} question:
\emph{on real production pipelines, does it work---and what decides whether it
does?} A static analyzer is only as good as the syntactic surface it recognizes, and
production data engineering is written by teams with internal frameworks, helper
libraries, and configuration-driven table names. The capability question can be
answered on benchmarks; the applicability question cannot.

\subsection{Why existing tools do not fill this gap}
Four bodies of prior work each address part of the problem but leave the core gap
open. \textbf{Runtime lineage tools} (Atlas~\cite{atlas}, OpenLineage~\cite{ol},
Spline~\cite{spline}, Unity Catalog) observe lineage from a running Spark driver;
they require the pipeline to be deployable and executable. \textbf{Code-semantic
graphs for LLMs} (CodexGraph~\cite{codexgraph}, GraphRAG~\cite{graphrag}) index
modules, classes, functions, and call edges---powerful for refactoring and
navigation, but silent on \emph{which columns flow into which derived columns}; a
\code{USES} edge does not distinguish control flow from data flow. \textbf{LLM-based
lineage extraction} (Yin et al.~\cite{yin}) prompts a model with whole pipeline
scripts, inheriting a per-script cost, latency, and non-determinism that AST
analysis avoids. \textbf{Static SQL lineage and impact analysis} is mature and
productized---\code{sqlglot}~\cite{sqlglot}, SQLLineage~\cite{sqllineage},
SQLMesh~\cite{sqlmesh}, and LineageX~\cite{lineagex} extract column lineage from a
SQL AST and use it for breaking-change analysis in CI---but every one of them
parses SQL (or dbt models compiling to SQL), \emph{not} the PySpark DataFrame API.
Finally, the only static PySpark analyzer, \code{pyspark-ast-lineage}~\cite{pyspark},
produces a flat set of table names: intermediate DataFrames are untracked,
transformations are not extracted, and no column-level information is emitted.

The gap is therefore specific: the static, column-level lineage that SQL tools
provide for SQL, and that runtime tools provide for Spark only by executing it, does
not exist for the \textbf{PySpark DataFrame API} at the source level---nor,
consequently, do the execution-free change-impact and version-diff capabilities it
would enable.

\subsection{From capability to applicability: what we found}
\label{sec:whatwefound}
We built the analyzer, and on benchmarks it works: 98\% semantic precision/recall on
hand-labeled lineage, and change-impact at 100\% precision and 78\% recall over 21
scenarios spanning synthetic pipelines and real third-party Databricks
Declarative-Pipeline code. Had we stopped there, the paper would report a capability
and leave its applicability unexamined---which is what we think the more interesting
question is.

We therefore evaluated on two proprietary production pipelines at an industrial
partner, using one for development and holding the second out behind a frozen tool
version. Three findings emerged, and they reshape the paper's claims.

\textbf{First, recovery is gated by syntax, not semantics.} The analyzer initially
recovered \emph{zero} column-level lineage on production code. Reaching 68\% impact
recall required handling nine concrete idioms---custom sink classes, dict-returning
table loaders, table names built as \code{f"\{database\}.\allowbreak silver.\allowbreak claim"}, bulk renames
via \code{toDF}, splatted and non-splatted projection lists---not one of which was a
limitation of the graph model, the impact algorithm, or the provenance resolver. Each
was an ordinary operation written in a form the extractor did not recognize. We
conclude that \emph{the barrier to static column lineage in production is surface-form
coverage, not transformation semantics}, and we contribute the inventory as a
starting point for anyone building such a tool.

\textbf{Second, applicability is domain-dependent, and our generalization claim does
not hold.} On the held-out pipeline---same organization, same medallion
architecture, same internal framework---impact accuracy was 0.167 precision and
0.118 recall, against 1.00 and 0.68 on the development pipeline. The difference is
structural: renames routed through a locally-defined function, and a gold source
table read from a runtime job parameter. The latter is beyond the reach of
\emph{any} static analyzer, and we distinguish such theoretic limits of
execution-free analysis from limitations of this implementation throughout.

\textbf{Third, the held-out corpus invalidated a claim rather than merely lowering a
score.} Our headline property---that impact analysis never reports a dependency it
did not recover---held on every benchmark and on the development pipeline, and broke
on the held-out one: an unresolved seed column silently fell back to matching that
column's \emph{name} in a different table, producing predictions that point backwards
along the dataflow. The defect was conditional on incomplete extraction, so no
quantity of additional development-set scenarios could have exposed it. We repair it,
verify that every previously reported number is unchanged, and state the precision
property conditionally from here on: DEGraph reports no false positives \emph{when it
has resolved the seed column's lineage}, and reports nothing when it has not.

We report the held-out number we obtained rather than the one we hoped for. For a
pre-merge gate, under-reporting is tolerable and confidently naming the wrong columns
is not; a tool of this kind should fail silent, and ours now does.

\subsection{Contributions}
\begin{enumerate}
\item \textbf{The DEGraph schema} --- a typed graph for data-engineering concepts.
Global nodes (\code{Table}, \code{ExternalSource}) represent persisted datasets
shared across files; file-local nodes (\code{DataFrame}, \code{Expression})
represent intermediate computation. Eight edge types encode column-level pipeline
structure, with provenance for conditional rules (\code{rule_logic}), array-struct
field lineage (\code{struct_fields}), and shared predicate dictionaries.
\item \textbf{A static, execution-free extractor} --- parses entire Databricks
repositories (\code{.py}, notebook exports, \code{.ipynb}, \code{.sql}) via Python's
\code{ast} module and \code{sqlglot}; no Spark cluster, no execution, no LLM at
extraction time. Accuracy against hand-labeled ground truth: \textbf{98\%
precision/recall/F1} semantic, 79\% exact-name. Against \code{pyspark-ast-lineage},
DEGraph uniquely recovers transformation edges, column provenance, and config-driven
fully-qualified table names.
\item \textbf{Static change-impact analysis for the PySpark DataFrame API} --- a
forward reachability query over inverted column provenance that returns every
downstream column/table affected by a change, deterministically and without
execution. On 21 scenarios (synthetic and real Databricks DLT code): \textbf{100\%
precision, 78\% recall, F1 88\%}; zero hallucinated dependencies \emph{when the seed
column's lineage is resolved}, and no output at all when it is not. Established for
SQL; we bring it to the PySpark DataFrame API.
\item \textbf{An industrial study of what breaks static PySpark lineage} --- two
proprietary production pipelines at an industrial partner, one held out behind a frozen tool
version. We contribute an inventory of \textbf{nine syntactic idioms} that determine
whether column lineage is recovered at all, evidence that the barrier is
\emph{surface-form coverage rather than transformation semantics}, a taxonomy
separating tool limitations from \emph{theoretic limits of execution-free analysis}
(a table name read from a runtime job parameter is resolvable by no static analyzer),
and a negative generalization result: 100\%/68\% on the development pipeline versus
0.167/0.118 on the held-out one. The held-out corpus also exposed a precision defect
that the development corpus structurally could not.
\item \textbf{Lineage version-diff with breaking-change detection} --- compares two
source revisions and reports added/removed tables, columns, and edges, plus the
downstream blast radius of removed/renamed columns; measured at 100\%
precision/recall on structural diff and correct breaking-vs-safe classification on
all six test edits, with zero false alarms.
\item \textbf{An honest LLM-context study (negative result)} --- a token-matched
BM25 baseline matched or beat the graph serialization for question-answering across
four model families; we retract the efficiency claim and report it as a cautionary
finding.
\end{enumerate}

\section{Related Work}
DEGraph sits among six threads; we state the gap narrowly to avoid overclaiming.

\paragraph{Static program analysis frameworks.} DEGraph is a static analysis, and it is
worth stating plainly where it sits with respect to that tradition. General frameworks for
analysing Python exist: MOPSA~\cite{mopsa} builds sound analyzers by abstract interpretation
and supports a large subset of Python~3; WALA~\cite{wala} provides a Python frontend; and
CodeQL~\cite{codeql} offers local and global data flow with taint tracking for Python. We do
not build on any of them, for two reasons.

The first is a difference in the soundness--precision tradeoff. A sound analysis
over-approximates: it models every behaviour the program might exhibit, and accepts false
positives as the price of missing nothing. DEGraph makes the opposite choice.
Change-impact is only actionable if an engineer can trust the answer, so the tool is built to
report nothing rather than guess, and its reported impact sets carry no false positives. An
over-approximating column analysis would return the columns a value \emph{might} reach, which
before a merge is close to useless. We arrived at this position empirically rather than by
design: the held-out study in \S\ref{sec:production} exposed a case where an unresolved seed
silently matched a same-named column in a different table, and the repair was to make
non-resolution explicit rather than to widen the search.

The second is that the gap is not Python semantics. All three frameworks analyse Python
values, variables and objects. A PySpark column is none of these: it is a string literal
passed into a method call on an opaque library object (\code{df.select("a")},
\code{df.withColumn("c", ...)}). Recovering column lineage therefore requires a domain
abstraction over the DataFrame API rather than a better analysis of the host language, and
that abstraction is where DEGraph's work sits. Our own central finding supports this
reading: the barrier we hit in production was surface-form coverage, not transformation
semantics (\S\ref{sec:production}), and a stronger general-purpose Python analyzer would not
have removed it.

The closest work in this tradition is Ariadne~\cite{ariadne}, which extends WALA with Python
support and a type system for tracking tensors, TensorFlow's core data structure, together
with a data flow analysis over their use. The shape of the problem is the same as ours: a
domain abstraction that the host language does not provide, layered onto program analysis.
Ariadne builds that abstraction inside a general framework; we build it directly over the
Python AST, trading generality for control over the DataFrame-specific surface forms that
\S\ref{sec:production} shows to be the binding constraint.

\paragraph{Runtime data lineage.} Atlas~\cite{atlas}, OpenLineage~\cite{ol} (with
Marquez~\cite{marquez}), and Spline~\cite{spline} capture lineage by observing
executing pipelines---attaching listeners to a Spark driver or dbt job. This rules
out code under review, branches that may never run, and lightweight LLM-assisted
workflows where provisioning a cluster is disproportionate. DEGraph operates
statically.

\paragraph{Code-semantic graphs for LLMs.} CodexGraph~\cite{codexgraph} indexes
Python repos into a graph of modules, classes, functions, and call edges for LLM
agents; GraphRAG~\cite{graphrag} and related systems follow the same principle.
These are powerful for call graphs and class hierarchies but silent on column-level
data flow. CodexGraph's own evaluation found its query-agent pattern consumed more
tokens than a BM25 baseline---a result that informs our delivery-mode framing.

\paragraph{LLM-based lineage extraction.} Yin et al.~\cite{yin} benchmark twelve
models that extract lineage tuples from whole pipeline scripts. The extractor is
itself an LLM call---with cost, latency, and stochasticity scaling per script---and
their corpus is Microsoft-internal and not reproducible. DEGraph extracts the same
column-level lineage deterministically via AST traversal, reproducibly, producing a
persistent artifact that supports execution-free impact and diff.

\paragraph{Static SQL lineage and change-impact.} The closest space, and most
important to position against: \code{sqlglot}~\cite{sqlglot} and
SQLLineage~\cite{sqllineage} extract column lineage from a SQL AST without
execution; SQLMesh~\cite{sqlmesh} uses it to flag breaking changes in CI; and
LineageX~\cite{lineagex} (ICDE 2025) does column-level lineage + impact from a SQL
AST. We do \emph{not} claim static column lineage or execution-free impact is novel:
for SQL it is established and productized. The gap is the \emph{language}---all of
these parse SQL (or dbt), none the PySpark DataFrame API.

\paragraph{Static PySpark analyzers.} \code{pyspark-ast-lineage}~\cite{pyspark}
parses Python ASTs to find read/write boundaries and resolve dynamic table names,
establishing feasibility, but emits a flat set of table names with no graph, no
transformations, and no column-level information.

\paragraph{Positioning.} What no prior system provides is the \emph{conjunction}:
static, source-AST, column-level lineage for the \textbf{PySpark DataFrame API},
serving as the substrate for execution-free impact and diff. DEGraph's contribution
is best understood as a \emph{technique transfer}---porting the proven SQL/dbt
static-lineage-plus-impact paradigm to the PySpark DataFrame API, with the new
extraction machinery (DataFrame variable tracking, transformation extractors,
declarative-pipeline and fluent-chain idiom support) that the language shift
requires.

\section{The DEGraph Schema}

\subsection{Design goals}
Six properties distinguish data-engineering lineage from general-purpose code
graphs: \textbf{repo-scale} (one unified graph in a single pass, no
inter-procedural analysis); \textbf{table-keyed} (files share data only through
named persisted tables, so \code{Table} names are the global cross-file join key);
\textbf{column-level} (track column identity through transformations);
\textbf{compactness} (a serialization that fits in a few thousand tokens);
\textbf{graceful degradation} (record an \code{<unresolved>} placeholder rather than
drop an edge); and \textbf{LLM-readability} (natural-key node identifiers).

\subsection{Two-layer node model}
\textbf{Global nodes} persist across files and form the cross-file backbone.
\code{Table} (a named persisted dataset: \code{fqn}, \code{format},
\code{columns[]}, \code{partition_cols[]}, \code{written_by}, \code{read_by[]}) is
deduplicated by fully-qualified name across the repo, so a write in one file and a
read in another wire to the same node. \code{ExternalSource} is a read-only source
not written within the repo. \textbf{File-local nodes} represent intermediate state
within one file: \code{DataFrame} (a PySpark variable or anonymous CTE) and
\code{Expression} (a derivation expression, stored once and referenced by id).

\subsection{Edge types}
Edges are typed, directional, and carry transformation metadata
(Table~\ref{tab:edges}). Three decisions warrant note. \textbf{Window functions
fold into \code{Derives}} (a \code{window_spec} field records partition/order/frame),
keeping the type set small. \textbf{\code{OpaqueTransform} prevents silent
disconnection}: when \code{df = helper(df)} cannot be analyzed, the opaque edge
preserves DataFrame identity rather than orphaning the downstream write.
\textbf{\code{dynamic=True} acknowledges enumeration failure}: a \code{*expr_list}
star-unpack built in a loop emits an edge with \code{<unresolved>} columns and a
note, keeping the graph connected.

Conditional rules and structured provenance are surfaced through four fields:
\code{DerivesEdge.\allowbreak rule_logic} (ordered \code{when/\allowbreak otherwise} chains, including those
embedded in \code{concat_ws}/\code{coalesce}); \code{Graph.\allowbreak column_rules}
(classification rules bound to Column-typed variables that never reach a
\code{withColumn} directly because they are passed into helpers);
\code{DerivesEdge.\allowbreak struct_fields} (per-field lineage of \code{array(struct(.\allowbreak .\allowbreak .\allowbreak ))}
columns, with NULL placeholders collapsed); and \code{Graph.\allowbreak predicate_dicts}
(shared dicts mapping category labels to predicate text).

\begin{table*}[t]
\centering
\caption{The eight DEGraph edge types.}
\label{tab:edges}
\footnotesize
\begin{tabular}{@{}L{2.4cm}L{5.9cm}L{6.7cm}@{}}
\toprule
Edge & Key payload & PySpark patterns matched \\
\midrule
\code{Reads} & \code{streaming}, \code{projected_cols[]} & \code{spark.table()}, \code{spark.\allowbreak read.\allowbreak .\allowbreak .\allowbreak load()}, SQL \code{FROM} \\
\code{Writes} & \code{mode}, \code{format}, \code{partition_cols[]} & \code{saveAsTable()}, \code{write.\allowbreak format().\allowbreak save()}, \code{writeStream.\allowbreak table()} \\
\code{Derives} & \code{output_col}, \code{source_cols[]}, \code{window_spec}, \code{rule_logic[]} & \code{withColumn}, \code{select}, \code{selectExpr}, SQL \code{SELECT expr AS col} \\
\code{Filters} & \code{predicate_id} & \code{filter()}, \code{where()}, SQL \code{WHERE} \\
\code{Projects} & \code{removed_cols[]} / \code{kept_cols[]} & \code{drop()}, column-set restriction \\
\code{Joins} & \code{join_type}, \code{join_keys[]}, \code{on_expr_id} & \code{join()}, \code{crossJoin()}, SQL \code{JOIN .\allowbreak .\allowbreak .\allowbreak  ON/\allowbreak USING} \\
\code{Aggregates} & \code{group_keys[]}, \code{agg_ops[]}, \code{output_cols[]}, \code{dynamic} & \code{groupBy().\allowbreak agg()}, SQL \code{GROUP BY} \\
\code{OpaqueTransform} & \code{operator}, \code{is_passthrough} & imported functions returning DataFrames \\
\bottomrule
\end{tabular}
\end{table*}

\subsection{Graph assembly}
A repo-level assembler discovers extractable files, sorts them (DDL/config first,
then bronze/silver/gold so source schemas are available), runs a \code{FileExtractor}
per file, and merges \code{Table} nodes by fully-qualified name (DDL column lists
take precedence; \code{written_by}/\code{read_by} are unioned). The key property is
that merging requires only string matching on table names---no inter-procedural
analysis, no call-graph traversal. Each file is independently correct; the global
graph emerges from the join. The schema is implemented as Pydantic~v2 models with
strict construction-time validation, so silent graph corruption surfaces as an
extraction warning rather than a wrong query answer.

\section{Implementation}

\paragraph{Input formats.} \code{.py} modules parse directly with \code{ast.parse};
notebook exports (with \code{\# COMMAND} cells and \code{\# MAGIC \%} directives) are
preprocessed to strip magics and resolve \code{\%run} imports; \code{.ipynb} cells
dispatch Python to the AST parser and SQL to \code{sqlglot}; \code{.sql} files parse
fully with \code{sqlglot} to populate the table-schema registry.

\paragraph{SafeEvaluator.} Dynamic table names
(\code{spark.\allowbreak table(f"\{database\}.\allowbreak silver.\allowbreak \{table\}")}) are resolved against a symbol
table populated from in-file literal assignments, \code{\%run}-imported assignments,
and \code{.degraph/\allowbreak helpers.\allowbreak json} overrides. It handles f-strings, concatenation,
attribute access, ternaries, and literal-key subscripts; unresolvable expressions
yield \code{<unresolved>} rather than failing.

\paragraph{DataFrame tracking and chain walking.} A forward-chain tracker maps
variable names to the active \code{DataFrame} node, creating a new node on each
reassignment (avoiding SSA while keeping pre/post states distinct). A chain-walker
unrolls nested method chains
(\code{df.\allowbreak filter(.\allowbreak .\allowbreak .\allowbreak ).\allowbreak drop(.\allowbreak .\allowbreak .\allowbreak ).\allowbreak withColumn(.\allowbreak .\allowbreak .\allowbreak )}) into ordered ops and groups
contiguous same-family operations, emitting one intermediate node per semantic
boundary rather than per method call.

\paragraph{Registries.} A helper registry maps function names to transformation
semantics (e.g.\ a \code{suffix_rename} helper emits one \code{Derives} edge per
renamed column); a sink registry maps custom sink classes (e.g.\ \code{UpsertSink})
to \code{Writes} edges with the right mode, so they are not silently skipped.

\paragraph{CTE parsing.} SQL embedded in \code{spark.\allowbreak sql(.\allowbreak .\allowbreak .\allowbreak )} is parsed by
\code{sqlglot} in Spark dialect; each named CTE in a \code{WITH} clause emits the
appropriate \code{Reads}/\code{Aggregates}/\code{Joins}/\code{Derives} edges, with
window functions populating \code{window_spec}.

\paragraph{Real-world Databricks idioms.} The imperative shape
(table-sourced DataFrame $\rightarrow$ variable-assigned chain $\rightarrow$
explicit write) dominates hand-written examples but not production code. An
external-validity study (\S\ref{sec:external}) drove support for three further
idioms, each routed back into the same chain-processing core: \emph{un-assigned
fluent pipelines} (a single \code{read}$\rightarrow$\code{transform}$\rightarrow$\code{write}
expression---fixed by descending the chain-walker through alternating method calls
and bare \code{.write}/\code{.writeStream} attributes); \emph{declarative pipelines}
(Delta Live Tables / Spark Declarative Pipelines, where a decorated function's
\emph{name} is the output table and its body \emph{returns} the pipeline); and
\emph{parameterized ingestion helpers} (inlined by binding call arguments to
parameters). A column-provenance refinement optimistically attributes a column of a
schema-less intermediate (e.g.\ a DLT table) to the table it was read from, while
tables with known schemas keep a strict membership check---leaving benchmark
provenance and impact precision unchanged.

\section{Evaluation}
We evaluate four claims. \S\ref{sec:extractor} measures how accurately the extractor
recovers lineage; \S\ref{sec:toolcmp} positions it against \code{pyspark-ast-lineage};
\S\ref{sec:impact} measures static change-impact and lineage diff (the
differentiator); \S\ref{sec:external} is the real-code external-validity study; and
\S\ref{sec:llm} reports the honest LLM-context negative result and a reliability
proxy. All extraction is static (Python AST + \code{sqlglot}); no pipeline is
executed and no LLM is invoked at extraction time.

\subsection{Benchmarks}
We use three hand-constructed PySpark repositories following the medallion pattern
(bronze ingestion $\rightarrow$ silver conformance $\rightarrow$ gold aggregation)
plus one real third-party corpus. Only \code{repo_synthetic_small} (13 files, 8
tables, 57 edges) has a full hand-labeled lineage graph; we use it as the accuracy
benchmark. The synthetic edge counts (57 / 80 / 163) are pinned by regression tests
so the numbers below remain reproducible. The real corpus is the
\code{databricks-demos/\allowbreak dbdemos}~\cite{dbdemos} retail churn pipeline (Spark Declarative Pipelines).

\subsection{Extractor accuracy}
\label{sec:extractor}
Recall alone is uninterpretable---an extractor that emits every edge scores 100\%
recall at near-zero precision. We measure both directions, per edge type, against
the hand-labeled graph, under two matching modes: \textbf{strict} (edge identity
includes intermediate DataFrame node names) and \textbf{semantic} (identity keys on
stable \code{table:}/\code{ext:} endpoints plus the per-edge content discriminator,
dropping volatile intermediate labels).

\begin{table}[h]
\centering
\caption{Extractor accuracy on \code{repo_synthetic_small} (57 edges).}
\label{tab:extractor}
\small
\begin{tabular}{@{}lccc@{}}
\toprule
Matching mode & Precision & Recall & F1 \\
\midrule
Strict (node-name-sensitive)   & 79\% & 79\% & \textbf{79\%} \\
Semantic (node-name-agnostic)  & 98\% & 98\% & \textbf{98\%} \\
\bottomrule
\end{tabular}
\end{table}

Under semantic matching, \code{reads}, \code{writes}, \code{derives},
\code{filters}, \code{joins}, \code{aggregates}, and \code{opaque_transform} all
match at 100\%; a single \code{projects} edge differs; table recovery is 8/8. The
79$\rightarrow$98 gap is a \textbf{labeling-convention artifact}: the ground truth
names each intermediate DataFrame with an idealized stage name
(\code{orders_trimmed} $\rightarrow$ \code{orders_normalized}), whereas the extractor
faithfully records the reused source variable (\code{orders = orders.\allowbreak withColumn(.\allowbreak .\allowbreak .\allowbreak )}
at every stage)---same operation, columns, and source columns, only a different
intermediate label. Semantic matching, the fair measure of lineage recovery,
confirms the structure is recovered correctly.

To test whether this is an artifact of self-authored code, we hand-labeled a second
ground-truth graph on a slice of \emph{real} Databricks code (the dbdemos retail
SDP transformations, 36 semantic edges, 7 tables), labeling from source alone. On a
blind first pass the extractor reached \textbf{100\% precision and 89\% recall (F1
94\%)}; the 100\% precision---zero false positives---confirms the no-hallucination
property holds on code we did not write. The recall gap was two general extractor
limitations (bare-\code{id} alias renames in a \code{select}; aggregates in a
DLT-function body), which we fixed; after the fixes, all 36 edges are recovered
(P/R/F1 = 100\%). We report the \textbf{89\% blind recall as the honest
generalization estimate}, since the fixes were informed by this ground truth.

\subsection{Comparison with pyspark-ast-lineage}
\label{sec:toolcmp}
We run both tools on \code{repo_synthetic_small}. To avoid a strawman we first
confirm the baseline is functional on its own bundled example (7 table/path strings
with detail records); a zero on our benchmark is therefore a real limitation, not a
setup failure. \code{pyspark-ast-lineage} produces a flat set of table/path strings
with no edges and no column lineage, and on our benchmark recovers \emph{nothing},
because the code uses idiomatic Databricks patterns---f-string fully-qualified names
resolved through a \code{\%run} config cell (\code{f"\{database\}.\allowbreak orders_raw"}) and
streaming reads---that defeat its variable resolver. DEGraph recovers all 8 tables
and 57 typed edges with column provenance. No existing static tool produces
column-level PySpark lineage, which is precisely the gap DEGraph fills.

\subsection{Static change-impact analysis and lineage diff}
\label{sec:impact}
\paragraph{Change-impact.} Given a column a developer is about to change, the
extractor's inverted column provenance yields a forward transitive closure---the set
of all downstream columns (and tables) affected, deterministically and without
execution. We evaluate against \textbf{21 ground-truth change scenarios}, each
authored from source by enumerating the genuinely-affected downstream columns. The
scenarios span all three synthetic benchmarks \emph{and} the real dbdemos DLT chain,
covering diverse edge types (derive, filter-referenced, aggregate, join/group key,
window) and hop depths, with the change-point at the bronze, silver, and gold tiers.
Real-code scenarios are evaluated against a committed extractor-output graph fixture
so they reproduce without the external repository.

\begin{table}[h]
\centering
\caption{Static change-impact accuracy (21 scenarios, micro-averaged).}
\label{tab:impact}
\small
\begin{tabular}{@{}lccc@{}}
\toprule
 & Precision & Recall & F1 \\
\midrule
\textbf{Pooled (21 scenarios)} & \textbf{100\%} & \textbf{78\%} & \textbf{88\%} \\
\quad \code{repo_synthetic_small} & 100\% & 94\% & 97\% \\
\quad \code{silver_clinical_claims} & 100\% & 86\% & 92\% \\
\quad \code{repo_synthetic_medium} & 100\% & 57\% & 73\% \\
\quad dbdemos retail (real DLT) & 100\% & 50\% & 67\% \\
\bottomrule
\end{tabular}
\end{table}

The \textbf{100\% precision (zero false positives, tp=49 fp=0 fn=14)} follows from
the construction: impact is graph reachability over edges the extractor actually
emitted, so a reported dependency always corresponds to a recovered one---a property
neither an LLM analyzer nor keyword retrieval can guarantee---and it holds on real
code we did not write, including the flagship dbdemos case where changing silver
\code{churn_users.\allowbreak creation_date} correctly reports the downstream gold feature
\code{churn_features.\allowbreak days_since_creation}. This property is \emph{conditional on the
seed column's lineage having been resolved}, and we state it that way deliberately:
\S\ref{sec:precision-defect} reports a held-out corpus on which unresolved seeds
caused the analyzer to fall back to a name match in a different table, violating the
property until we repaired it. Every number in this section is unchanged by that
repair. The \textbf{78\% recall} reflects a
resolver that now threads column lineage through \code{spark.sql} blocks: we extended
the CTE handler so every computed column in a CTE \emph{or} an outer SELECT (e.g.\
\code{inv.\allowbreak on_hand /\allowbreak  NULLIF(ds.\allowbreak sold,0) AS turnover_rate}) carries the columns its
expression references, lifting recall from 67\% to 78\% at unchanged 100\% precision.
The residual misses are a small, named set: (i) Python-side \code{coalesce}/\code{when}
derivations applied \emph{after} a SQL block, (ii) group-key$\rightarrow$aggregate-value
dependencies (a \code{GROUP BY} key affecting the aggregate it buckets), and (iii)
DLT-function-body subquery aggregates. Ground truth follows a value-dependent
definition of impact (a column is affected iff its value would change), enumerated
from source; an earlier resolver round surfaced and we corrected two labeling errors
(one phantom intermediate-column entry, several omissions on the highest-fan-out
scenario), documented in \code{experiments/\allowbreak impact_eval.\allowbreak py}.

\paragraph{Lineage version-diff.} The differ compares two source revisions and
reports added/removed tables, columns, and edges, plus the downstream \textbf{blast
radius} of removed/renamed columns (computed on the old graph---what the change
breaks). We evaluate on six realistic edits---three breaking (renaming a silver
aggregate or timestamp that gold still consumes) and three safe (dropping/renaming a
leaf column, adding a new column)---re-extracting both revisions
(Table~\ref{tab:diff}).

\begin{table}[h]
\centering
\caption{Lineage version-diff accuracy (six edits).}
\label{tab:diff}
\small
\begin{tabular}{@{}ll@{}}
\toprule
Diff capability & Result \\
\midrule
Structural diff (added/removed cols) & P 100\% / R 100\% / F1 100\% \\
Breaking-change classification & 3/3 breaks caught, 3/3 safe not alarmed \\
Blast-radius recall (breaking edits) & 85\% (inherits impact recall) \\
\bottomrule
\end{tabular}
\end{table}

The structural diff is exact and the breaking-vs-safe verdict is correct on every
edit: the differ catches all three real breaks (e.g.\ renaming silver
\code{customer_profile.\allowbreak lifetime_revenue} $\rightarrow$ \code{ltv_total} is flagged
breaking because gold \code{customer_ltv} still references it) while raising
\textbf{no false alarms} on the three safe edits---the property a pre-merge gate
needs to be trusted. Blast-radius recall (85\%) inherits the same residual
limitation as change-impact. This is the pre-merge / CI use case that runtime tools
(needing execution) and retrieval (lacking structure) cannot serve.

\subsection{External validity: from public code to production pipelines}
\label{sec:external}
\label{sec:production}
The benchmarks above are author-written or public tutorial code. We now report what
happened when the extractor met code written by other people---first a public
third-party corpus, then two proprietary production pipelines. The two studies tell
the same story at different scales, and we present them together because that
repetition is the finding.

\paragraph{A first signal from public code.} Running the extractor on
\code{dbdemos-notebooks} (retail churn) and \code{pyspark-examples} produced a
negative initial result: files parsed without crashing but recovered almost no
transformation lineage---one derive across 17 dbdemos files. The cause was not
analytical depth but syntactic reach: the extractor recovered lineage only for the
imperative table$\rightarrow$variable-chain$\rightarrow$explicit-write idiom of our
synthetic benchmarks, while the three idioms dominating real Databricks code
(Declarative-Pipeline-decorated functions, un-assigned fluent chains, parameterized
ingestion helpers) each produced \emph{zero} edges. Adding support for those three
(\S4) took dbdemos retail from 12 to 33 to \textbf{64} edges---a $5.3\times$
increase---\emph{with zero change to the synthetic graphs} (57/80/163, pinned by
regression tests), and a precision spot-check confirmed the new edges are correct.
With lineage recovered, both impact and diff then ran on that corpus: changing
silver \code{churn_users.\allowbreak creation_date} correctly reports downstream impact on
\code{churn_features}, which is, to our knowledge, the first static, execution-free,
column-level change-impact analysis demonstrated on real Databricks
Declarative-Pipeline code.

\paragraph{Two production pipelines at an industrial partner.} To test whether that pattern holds
against production engineering rather than public demonstration code, we evaluated
DEGraph on two proprietary Databricks medallion pipelines at an industrial partner:
\textbf{Pipeline A} (development) and \textbf{Pipeline B} (\textbf{held
out}). The corpora cannot be redistributed. Every number below is reproducible from
released \emph{anonymized} bundles containing the hashed graph, the hashed
scenarios, and the scores: names are replaced by consistent salted hashes, and
precision/recall/F1 are set-intersection quantities, hence invariant under that
bijection---each run asserts the anonymized scores equal the real ones before
writing the bundle. Change-impact labels were drafted by an LLM reading source only
and then verified and corrected by the first author against the code; we report this
rather than describing them as hand-labeled.

\paragraph{Recovery requires surface-form coverage, not deeper semantics.}
DEGraph initially recovered \emph{zero} column-level lineage on Pipeline A
---the same failure mode as dbdemos, one order of magnitude larger. Reaching usable
coverage required handling nine further idioms (Table~\ref{tab:idioms}). As with the
public corpus, none was a deficiency of the graph model, the impact algorithm, or the
provenance walk: in every case the pipeline expressed an ordinary operation in a
syntactic form the extractor did not recognize. Two were pure configuration; seven
were extractor changes, each developed against a synthetic reproduction of the
idiom---never against the proprietary source---each covered by a regression test, and
each verified to leave the synthetic benchmark bit-identical. Having now observed the
pattern on two independent corpora, we state it as the study's central empirical
claim: \emph{the barrier to static column lineage on production PySpark is
surface-form coverage, not transformation semantics}.

\begin{table}[h]
\centering
\caption{Idioms that had to be handled before any lineage was recovered.}
\label{tab:idioms}
\small
\begin{tabular}{@{}L{7.6cm}cL{6.0cm}@{}}
\toprule
Idiom & Fix & Effect if absent \\
\midrule
Custom sink classes, \code{UpsertSink(\ldots).\allowbreak save(df)} & config & no writes at all \\
Dict-returning table loader, \code{tbls["t"]} & code & no schema anchor \\
\code{union} / \code{unionByName} chains & code & tracking severed \\
Helper with the DataFrame in a non-first argument & code & lineage absorbed \\
Splatted projection, \code{select(*cols)} & code & result orphaned \\
Config-driven table names, \code{f"\{db\}.\allowbreak silver.\allowbreak t"} & code & one table becomes 2--3 nodes \\
Bulk rename, \code{toDF(*[c.\allowbreak lower() for c in df.\allowbreak columns])} & code & snaps after the read \\
Non-splatted projection, \code{select(cols)} & code & result orphaned \\
Seed resolution crossing tables & code & \emph{precision} loss (\S\ref{sec:precision-defect}) \\
\bottomrule
\end{tabular}
\end{table}

\paragraph{Development pipeline (Pipeline A).} Twenty-nine files, 84 tables, 2{,}808
edges. On 15 change scenarios, impact accuracy improved across successive idiom
fixes from 21\% to 44\% to \textbf{68\% recall}, and \textbf{precision was exactly
100\% at every stage}---zero false positives in all 15 scenarios, across eight fixes
and a $3.2\times$ change in recall, with 8 of 15 scenarios at perfect recall. Because
three of those fixes were developed against failures this set exposed, we report it
as a development trajectory, not a generalization estimate. Of the 24 residual
misses, 21 are composite-column constructors (\code{struct}, \code{array},
\code{map\_from\_entries}, \code{collect\_list}): DEGraph models column-to-column
derivation and does not descend into composites.

\paragraph{Held-out pipeline (Pipeline B).} We tagged and froze the extractor before
touching this corpus, restricted adaptation to configuration, declared the sampling
rule in advance (five bronze$\rightarrow$silver, four silver$\rightarrow$silver,
three silver$\rightarrow$gold), and evaluated once. Pooled accuracy was
\textbf{precision 0.167, recall 0.118, F1 0.138} (tp=2, fp=10, fn=15); counting three
false positives that are genuine consumers our labels failed to name gives an upper
bound of 0.417/0.250/0.313. Recovery was confined to bronze$\rightarrow$silver
(P 0.667, R 0.222); silver$\rightarrow$silver and silver$\rightarrow$gold recovered
nothing. The 15 misses divide into three classes: eight \emph{tool limitations}
(seven from a bulk rename routed through a locally-defined function---an operation
DEGraph handles inline but does not follow one level into a local \code{def}), three
composite-column bundles, and four that are \emph{statically undeterminable} because
the gold source table is a runtime job parameter supplied by
\code{dbutils.\allowbreak widgets.\allowbreak get}. That last class is a limit of execution-free analysis
itself, not of this implementation: no static analyzer can resolve it.

\paragraph{The generalization claim, tested properly, does not hold.} The gap between
1.00/0.68/0.81 on Pipeline A and 0.17/0.12/0.14 on Pipeline B---two pipelines from the same
organization, sharing a medallion architecture and a shared internal framework---is
structural rather than sampling noise. DEGraph recovers lineage where pipelines map
columns directly and goes nearly dark where they route renames through local helper
functions and take table names from runtime parameters. We therefore make no claim of
uniform applicability; the scope condition is now characterized rather than unknown.

\subsection{A precision defect that only a held-out corpus could expose}
\label{sec:precision-defect}
The held-out run did not merely lower a score; it invalidated a claim. Seven of the
ten false positives share one signature. When a seed column had no resolved lineage
of its own, seed resolution fell back to matching the column \emph{name} across every
table and returning that table's descendants. Because column names repeat across a
medallion pipeline, this produced predictions pointing \emph{backwards} along the
dataflow: seeding \code{silver.\allowbreak customer.\allowbreak audit\_record\_id} returned
\code{silver.\allowbreak wh\_customer.\allowbreak *}, although \code{wh\_customer} feeds \code{customer} and not
the reverse. The same fallback produced the run's few accidentally-correct
predictions, so the tool was not reasoning about those either.

This defect was structurally invisible to every prior evaluation. Where lineage
extracts cleanly, seeds resolve through real edges and the fallback never executes;
no quantity of additional development-set scenarios could have surfaced it. We
repaired seed resolution so that the table constraint is never dropped: match the
qualified name, else the same table by short name, else an unqualified file-local
column node only when no other table qualifies that column anywhere. An ablation over
five candidate rules confirms the diagnosis---the cross-table path contributed
\emph{nothing} to the Pipeline A result and produced every real false positive on Pipeline B.
The repair costs one true positive on Pipeline A (recall $0.680 \rightarrow 0.667$,
precision unchanged) and removes all ten false positives on Pipeline B. We re-ran every
previously reported result: the 21-scenario impact evaluation
(Table~\ref{tab:impact}) and the diff evaluation (Table~\ref{tab:diff}) are
bit-identical before and after.

We draw two conclusions. First, on the substance: \emph{silent under-reporting is
tolerable in an impact-analysis tool; confidently naming wrong columns is not}, so a
seed the analyzer cannot resolve must yield nothing rather than a guess. Second, on
method: this is a concrete demonstration that a held-out corpus buys something a
larger development corpus cannot, because the defect was conditional on the very
incompleteness that the development corpus lacked.

\subsection{LLM-context study: an honest negative result, and a reliability proxy}
\label{sec:llm}
The project's original thesis was that a compact graph serialization is a superior
\emph{LLM context} format. We tested this against raw source and a token-matched
BM25 retrieval baseline~\cite{bm25}, across four model families and three benchmarks, on a
sealed test split evaluated once per frozen tool tag.

\begin{table}[h]
\centering
\caption{LLM-context Q\&A (Gemini, pooled): graph serialization vs.\ baselines.}
\label{tab:llm}
\small
\begin{tabular}{@{}lcc@{}}
\toprule
Context & Accuracy & Tokens \\
\midrule
Raw source & 85.8\% & baseline \\
BM25 retrieval (token-matched) & 81.7\% & $\approx$ graph \\
Compact graph serialization & 75.0\% & $\approx$ BM25 \\
\bottomrule
\end{tabular}
\end{table}

\textbf{At equal token budget, a trivial lexical baseline matched or beat the graph
serialization.} We therefore retract the efficiency claim and report this as a
cautionary finding for the ``graphs-as-LLM-context'' literature, which under-reports
the strength of BM25. DEGraph's value is the deterministic, execution-free
structural analysis of \S\ref{sec:impact}, not token reduction.

\paragraph{Reliability proxy for the impact ground truth.} Because all ground truth
is single-author, we add an automated reliability proxy (a human second annotator
with Cohen's $\kappa$ remains future work). Two independent LLM judges from different
model families (GPT-4o, Gemini 2.5 Flash) were shown the pipeline source and each
change, then a shuffled candidate list mixing the ground-truth-affected columns with
sampled non-affected columns---\emph{unlabeled}---and asked to select the affected
ones from source alone. Over 229 rated items per judge, agreement with the author's
labels was \textbf{Cohen's $\kappa=0.86$ (GPT-4o) and $0.91$ (Gemini 2.5 Flash)}
---``almost perfect'' on the Landis--Koch scale---at 94--97\% raw agreement. This is
an LLM-as-judge proxy, \emph{not} human inter-rater agreement. The disagreements are
directional: judges almost never miss a labeled dependency (author-positive,
judge-negative = 3 and 0 of 61) but occasionally add one the author conservatively
excluded (join/partition-key columns, where a judge takes a broad ``the whole join
changes'' view). The ground-truth positives are thus strongly corroborated and the
labels are, if anything, conservative.

\section{Discussion}

\subsection{Threats to validity}
\textbf{Single annotator, mitigated by a proxy.} All ground truth was authored by
one annotator. A second \emph{human} annotator with Cohen's $\kappa$ is the gold
standard and top future work; in its place we report the LLM-judge proxy above
($\kappa=0.86$/$0.91$). \textbf{Public-corpus external validity by coverage + spot-check.} Public real
corpora have no hand-labeled lineage, so the dbdemos portion of
\S\ref{sec:external} reports coverage and spot-checked precision rather than full
P/R, except on the one hand-labeled real slice (\S\ref{sec:extractor}).
\textbf{Statistical power of the LLM study.} The Q\&A sets (10--25 questions each)
are below the size for formal significance testing; the negative result is reported
as directionally consistent across four families and three benchmarks, and the
contributions do not depend on it.

\textbf{Development-set contamination, and what we did about it.} Pipeline A
is \emph{not} an unbiased estimate: three extractor fixes were developed
against failures it exposed, so we report it as a trajectory and never as a
generalization number. To obtain an unbiased estimate we tagged and froze the
extractor before Pipeline B was opened, restricted adaptation to
configuration (sink and helper registries, catalog schemas) with no code changes,
declared the scenario-sampling rule in advance, and evaluated once. During that run
we were advised to re-scope the scenarios away from layer transitions already known
to fail; we declined, since choosing the test after seeing which paths work voids
the estimate as surely as tuning does. The repair prompted by that run
(\S\ref{sec:precision-defect}) was adopted only after verifying it left every
previously reported result bit-identical and cost exactly one true positive on the
development set. The Pipeline B result is nonetheless a \emph{one-shot} estimate on a
single corpus, and its confidence interval is correspondingly wide.

\textbf{Construct validity of the impact labels.} Ground truth follows a
value-dependent definition---a column is affected iff its value would
change---which counts pure passthroughs as affected and excludes join and partition
keys that influence row matching but not values. A reader adopting a
row-membership definition would obtain different sets. Labels were LLM-drafted from
source and author-verified; on Pipeline B this process produced two of its
own errors (seed names recorded in post-transformation casing), which we found while
classifying misses, corrected, and disclose along with what did \emph{not} change.
Three further member false positives proved to be genuine downstream consumers our
labels failed to name, so we also report an upper bound treating them as correct.

\textbf{Generality of the idiom inventory.} Our production evidence is two pipelines
from a \emph{single} organization sharing one internal framework, evaluated by one
author. The nine idioms should be read as a lower bound on the diversity a
multi-organization study would find, and the observed gap between the two pipelines
as evidence that such diversity exists \emph{within} an organization, not as a
measurement of how DEGraph performs on production PySpark in general.

\subsection{Limitations}
\textbf{Where DEGraph is deliberately incomplete, and why we say so.} The limitations below
are not a list of defects awaiting repair. Most are consequences of two design choices made
knowingly: the analysis is \emph{intra-procedural}, resolving lineage within a file rather
than across function boundaries, and it does not perform \emph{alias analysis}, so two names
bound to the same DataFrame are not identified as one. Both choices buy the precision
property this paper leans on, and both cost recall in exactly the places our held-out study
found. We state them in these terms because a static analysis that is unsound in specific,
known ways should say which ways: the alternative, as the soundiness
manifesto~\cite{soundiness} puts it, is that sources of unsoundness ``lurk in the shadows,
with caveats only mentioned in an off-hand manner in an implementation or evaluation
section.'' The held-out pipeline in \S\ref{sec:production}, run behind a frozen extractor, is
our attempt to measure the extent of that unsoundness rather than assert its absence.

\textbf{Impact recall (78\%) is bounded by three small, named gaps} after the
\code{spark.sql} CTE/outer-SELECT column-lineage resolver was extended (which lifted
recall from 67\% to 78\% at unchanged 100\% precision): (i) Python-side
\code{coalesce}/\code{when} derivations applied after a SQL block; (ii)
group-key$\rightarrow$aggregate-value dependencies; and (iii) DLT-function-body
subquery aggregates. Precision is unaffected throughout (the tool never invents a
dependency); each gap is localized rather than systemic.
\textbf{Real-code diff} column-detection is bounded because DLT
\code{select(.\allowbreak .\allowbreak .\allowbreak ).\allowbreak alias()} provenance keys on the source column rather than the
output alias; the real-code diff is demonstrated with a direct downstream query
rather than the automatic breaking-list. \textbf{Two-statement inter-procedural
ingestion helpers} (separate read and write statements inside a function) are not
yet inlined.

The production study (\S\ref{sec:production}) adds four limitations we consider more
consequential than the above. \textbf{Interprocedural renames:} a bulk rename routed
through a locally-defined function is not followed, though DEGraph handles the same
operation inline; this alone accounts for seven of fifteen held-out misses and is the
single most valuable fix available. This is the intra-procedural boundary stated concretely:
nearly half of the held-out misses are one design choice, not a scatter of unrelated bugs,
which is what makes the fix tractable rather than open-ended. \textbf{Composite-column provenance:} DEGraph
does not descend into \code{struct}, \code{array}, \code{map\_from\_entries}, or
\code{collect\_list} constructors, making this the largest miss class in both
production corpora. \textbf{Runtime-parameterized table names:} when a table is named
by a job widget, no static analysis can recover it---we report such cases as
undeterminable rather than as tool error. \textbf{Case sensitivity:} Spark resolves
identifiers case-insensitively while DEGraph matches them case-sensitively, so a
table referenced with inconsistent casing yields duplicate nodes. Finally, our
production evidence is two pipelines from a \emph{single} organization sharing one
internal framework; the idiom inventory should be read as a lower bound on the
diversity a broader study would find.

\section{Conclusion}
The answer to our central research question is affirmative but \emph{conditional}:
static analysis of PySpark source can produce a column-level lineage graph accurate
enough to support deterministic, execution-free change-impact analysis and lineage
version-diff---on synthetic benchmarks, on real third-party Databricks code, and on a
production pipeline at an industrial partner---but whether it does so on any \emph{particular}
pipeline depends on how that pipeline is written. DEGraph occupies a precise gap
that SQL tools (which cannot parse DataFrame code), runtime Spark tools (which
require execution), and the one existing static PySpark analyzer (which stops at
table names) leave open. Within that gap, its accuracy is bounded less by the
analysis than by the syntactic surface it must recognize.

Our industrial evaluation is the part we expect to be most useful to others. Nine
concrete idioms stood between zero and 68\% impact recall on a real production pipeline,
and none of them was a limitation of the graph model. On a second pipeline from the
same organization, held out behind a frozen tool version, impact recall was 0.118 and
precision 0.167: renames routed through a local helper function and a table name read
from a runtime job parameter put most of that pipeline beyond reach---the latter
beyond the reach of any static analyzer. That run also exposed a precision defect
invisible to every prior evaluation, because it fired only where extraction was
incomplete. We report the number we obtained, not the one we hoped for; a tool of
this kind must under-report rather than name the wrong columns, and it now does.

The highest-value future work follows directly: interprocedural analysis, so a rename
inside a local function is followed; provenance \emph{into} composite constructors,
the largest single miss class in both corpora; a multi-organization study to
establish how the idiom inventory saturates; and an independent second \emph{human}
annotator to complement the automated reliability proxy ($\kappa=0.86$/$0.91$).

\paragraph{Reproducibility.} The extractor, the three synthetic benchmarks with
hand-labeled ground truth, the impact and diff evaluations, the committed real-code
graph fixture, and the LLM-judge harness are released with this paper.

\begin{thebibliography}{99}
\small
\bibitem{atlas} Apache Atlas, 2024. Apache Software Foundation. \url{https://atlas.apache.org/}
\bibitem{ol} OpenLineage, 2024. The Linux Foundation. \url{https://openlineage.io/}
\bibitem{marquez} The Marquez Project, 2024. \url{https://marquezproject.ai/}
\bibitem{spline} AbsaOSS, 2024. Spline: Data lineage tracking and visualization for Apache Spark. \url{https://absaoss.github.io/spline/}
\bibitem{codexgraph} X. Liu et al., 2025. CodexGraph: Bridging Large Language Models and Code Repositories via Code Graph Databases. \emph{Proc.\ NAACL 2025}. arXiv:2408.03910.
\bibitem{graphrag} D. Edge et al., 2024. From Local to Global: A Graph RAG Approach to Query-Focused Summarization. Microsoft Research.
\bibitem{yin} J. Yin, Y.-W. Chen, M.-L. Lee, X. Liu, 2025. Schema Lineage Extraction at Scale: Multilingual Pipelines, Composite Evaluation, and Language-Model Benchmarks. arXiv:2508.07179.
\bibitem{sqlglot} T. Mao, 2024. SQLGlot: a no-dependency SQL parser, transpiler, and optimizer with column-level lineage. \url{https://github.com/tobymao/sqlglot}.
\bibitem{sqllineage} Reata, 2024. SQLLineage: SQL Lineage Analysis Tool powered by Python. \url{https://github.com/reata/sqllineage}.
\bibitem{sqlmesh} Tobiko Data, 2024. SQLMesh: column-level lineage and change-impact analysis for SQL data pipelines. \url{https://sqlmesh.com/}.
\bibitem{lineagex} S. Zhang et al., 2025. LineageX: A Column-Level Lineage Extraction System for SQL. \emph{Proc.\ IEEE ICDE 2025} (demo). arXiv:2505.23133.
\bibitem{pyspark} R. Espinosa, 2024. pyspark-ast-lineage. \url{https://github.com/richardesp/pyspark-ast-lineage}.
\bibitem{ariadne} J. Dolby, A. Shinnar, A. Allain, J. Reinen, 2018. Ariadne: Analysis for Machine Learning Programs. \emph{Proc.\ 2nd ACM SIGPLAN Int.\ Workshop on Machine Learning and Programming Languages (MAPL)}, 1--10. arXiv:1805.04058.
\bibitem{mopsa} M. Journault, A. Min\'e, R. Monat, A. Ouadjaout, 2019. Combinations of Reusable Abstract Domains for a Multilingual Static Analyzer. \emph{Proc.\ VSTTE 2019}. \url{https://mopsa.lip6.fr/}
\bibitem{wala} IBM Research, 2024. WALA: T.J.\ Watson Libraries for Analysis. \url{https://wala.github.io/}
\bibitem{codeql} GitHub, 2024. CodeQL: data flow and taint tracking libraries for Python. \url{https://codeql.github.com/docs/codeql-language-guides/analyzing-data-flow-in-python/}
\bibitem{soundiness} B. Livshits, M. Sridharan, Y. Smaragdakis, O. Lhot\'ak, J. N. Amaral, B.-Y. E. Chang, S. Guyer, U. Khedker, A. M\o{}ller, D. Vardoulakis, 2015. In Defense of Soundiness: A Manifesto. \emph{Commun.\ ACM} 58(2), 44--46.
\bibitem{bm25} S. Robertson, H. Zaragoza, 2009. The Probabilistic Relevance Framework: BM25 and Beyond. \emph{Foundations and Trends in IR} 3(4), 333--389.
\bibitem{dbdemos} Databricks, 2024. dbdemos: Databricks demos and accelerators. \url{https://github.com/databricks-demos/dbdemos-notebooks}.
\end{thebibliography}

\end{document}
